\chapter{Conclusion}
\label{chap:conclusion}

This thesis developed and evaluated imputation methods for regression inference in two-phase validation designs. An error-prone measurement is observed for the full cohort, while a reference measurement is observed only in a small validated subsample. Probability designs specify inclusion probabilities, whereas deterministic tail designs specify a validated region and require extrapolation beyond it. When only five to ten per cent of the cohort is validated, most cohort members lack the reference measurement. Imputation methods complete those values; design-based methods instead use weighting or calibration. The target is the regression coefficient on the reference scale and its confidence interval. Individual prediction error is secondary to that inferential target. Chapter~\ref{chap:SimDesigns} therefore evaluates coefficient bias, relative root mean squared error, nominal coverage and two oracle coverages. The oracle quantities show how coverage changes after replacing either the reported standard error or the average simulation bias.

The thesis proposes three methods under a common comparison protocol \citep{amorim2021two}. Sample Imputation via Sample Likelihood (SISL) addresses informative sampling when a design weight is released but parts of the sampling information are hidden. Sample Imputation via Conditional Generative Adversarial Networks (SICG) adapts conditional tabular generation to the two-phase missing block. Sample Imputation via Conditional Residual Denoising Diffusion Networks (SIRD) uses the phase-one proxy as a reference endpoint and fits a separate denoiser for each bootstrap imputation. Their results depend on the information released to the analyst and on whether the target is a coefficient estimate or an interval.

\section{Two information settings for two-phase imputation}
\label{sec:conc_regimes}

The main division is between a design whose selection variables are partly hidden and designs whose selection-related phase-one variables are available to the imputation model.

Simulation Design~I forms 27 sampling strata from the observed outcome tertile and the hidden stratifying variables $C=(S_1,S_2)$. The analyst observes the outcome component of the joint stratum and the released design weight, but not $S_1$, $S_2$ or $C$. A naive model fitted to the validated rows has exposure bias of $13.4\%$. Generalised raking reduces it to $0.3\%$, weighted-likelihood multiple imputation reduces it to $0.9\%$, and SISL reduces it to $2.2\%$. SISL is therefore a negative main result. It removes most of the naive bias, but its residual $2.2\%$ exposure bias does not reach weighted-likelihood imputation or raking, and that residual is a property of the weight model rather than of the variance construction: with the stratifying variables suppressed, the estimated inclusion surface cannot exactly recover the true inclusion function (Chapter~\ref{chap:SISL}), so a small bias survives however the interval is formed. The three variants share that bias and differ only in the covariance used for the interval, where the sandwich and bootstrap versions reach near-nominal coverage and the inverse-information version under-covers. The result is specific to this design and these working models, and does not establish a general failure of sample-likelihood imputation under informative sampling.

Adding the factor-coded design weight to the naive imputation model did not help in the tested extension: exposure bias was $14.2\%$, compared with $13.4\%$ without it. Treating the weight as this predictor was not equivalent to the likelihood weighting used by weighted-likelihood multiple imputation.

In Simulation Designs~II and~III, the phase-one variables related to selection are available to the imputation model. Conditional identification still requires the independence and support conditions stated in Chapters~\ref{chap:SICG} and~\ref{chap:SIRD}; deterministic tail rows also test extrapolation. SICG gives low aggregate bias on the continuous block and reasonable categorical estimates, while its retained MC-dropout intervals are often too narrow. The complete SIRD procedure improves the continuous and categorical coefficient estimates and coverage relative to SICG. Its bootstrap fits supply fitted-model variation for the intervals; the study does not attribute the coefficient-estimate differences to bootstrapping. In Design~III its continuous nominal coverage is $0.89$ with an SE ratio of $1.37$, so the reported standard errors are conservative on average even though coverage remains below $0.95$. Its categorical nominal coverage is $0.99$ with an SE ratio of $1.40$. Generalised raking retains lower aggregate bias where its calibration constraints apply. Multiple imputation instead returns $m$ reusable completed datasets, which can support later analyses without constructing a new calibration estimator for every target \citep{Rubin1987,breslow2009improved,han2021combining}.

The continuous-block relative RMSE in Design~III is $2.5\%$ for SIRD, $2.9\%$ for SICG and $3.5\%$ for raking. In VCCC, the corresponding values are $6.1\%$, $7.1\%$ and $7.0\%$. SIRD has the lowest continuous-block relative RMSE among these three methods in both settings, while raking retains the smaller bias. On the Cox product term and in Design~I, the imputation methods have larger displacement. The relative advantage therefore depends on both the design information and the analysis target.

\section{Coefficient estimates and fitted-model uncertainty}
\label{sec:conc_uncertainty}

A single trained network can produce many completed datasets by varying its sampling noise. Those datasets do not vary the fitted network unless the procedure adds another source of fitted-model variation. Rubin's between-imputation term records variation among completed-data estimates. It includes fitted-model uncertainty only when the imputation procedure propagates that uncertainty. Low coefficient bias and low distributional discrepancy at one fitted parameter do not establish interval calibration.

Proposition~\ref{prop:improper-mi} gives a first-order variance comparison under conditions (C), (M) and (R). Under those conditions, fixing the fitted imputation model omits the parameter-uncertainty component represented in the proposition. Its scope is a first-order comparison, rather than an unconditional finite-sample coverage theorem. Lemma~\ref{lemma:mean-collapse} is also a limiting result: as a squared-error reconstruction penalty dominates the remaining objective, the generated values concentrate on the conditional mean. The finite SICG comparison then shows a $9.5\%$ reduction in marginal imputation spread and a $72\%$ reduction in $B$ under the tested reconstruction condition. Uncongeniality, a mismatch between the imputation and analysis models, is a separate issue again: depending on the relation between the two it can displace the completed-data estimate, as on the Cox age-by-insulin interaction of Section~\ref{sec:conc_limitations}, or leave Rubin's variance conservative or anti-conservative \citep{meng1994multiple,robins2000inference}.

The narrowing of the imputed spread is not confined to SICG. Figure~\ref{fig:appc_simI_dispersion} in Appendix~\ref{app:suppC_simI} repeats Figure~\ref{fig:gain_underdispersion} for each imputation method on one Design~II replicate, plotting the outcome against the exposure that the method placed in the unvalidated rows. A method that imputes from a single fitted model tends to pull the imputed exposure towards its conditional mean, so the imputed spread is narrower than the true spread and the fitted line in the panel is steeper than the true line. MIXGB shows this most strongly, with the imputed exposure held in a narrow band; GAIN and MICE-RF show it to a smaller degree. SICG and SIRD keep a spread near the reference, and their fitted lines stay near the true line. The narrowing is already present in the completed data, before the analysis model is fitted, and it is the same distortion that the Design~II coefficient bias reflects.

The same distributional view can be taken for the imputed outcome. Figure~\ref{fig:appc_simIII_TIdensity} in Appendix~\ref{app:suppC_simII} shows the imputed survival time under Simulation Design~III with outcome error, on the unvalidated units. The true survival time is right-skewed with a pile-up at the censoring bound, and the methods reproduce this shape with different success. SIRD and MICE-RF stay close to the true density. SICG is too concentrated near the mode, GAIN moves the mode to shorter times, MIXGB adds secondary modes that are not present in the truth, and MICE-PMM flattens the body while raising the spike at the censoring bound. This is a property of the completed values themselves, and it records which methods recover the shape of the imputed outcome. That is a separate question from whether the fitted coefficient and its interval are calibrated.

The simulations provide separate evidence about the implemented uncertainty procedures. SICG's continuous-block SE ratio is $0.32$ in Design~II and $0.59$ in Design~III. Within each diagnostic row, its tested MC-dropout procedure produces less $B$ than bootstrap retraining. SIRD uses the proxy-referenced residual construction with 20 bootstrap denoiser fits. Its continuous-block SE ratios are $1.13$ and $1.37$, with nominal coverages of $0.93$ and $0.89$. Bootstrap refitting raises the expected between-imputation variance above the single-fit baseline, but it overshoots: these two SE ratios, together with the categorical-block value of $1.40$, all sit above one, so the reported intervals are conservative rather than calibrated. Raising the between-imputation variance is not the same as scaling it correctly, and a correctly calibrated propagation of fitted-model uncertainty is improved but not achieved. In Chapter~\ref{chap:SISL}, the three variants share one point estimate and differ only in the covariances used for the redraws. SISL-M uses model-based covariances at both fitted stages and has coverage of $0.83$; the sandwich and bootstrap versions use robust and resampling covariances and reach $0.94$ and $0.96$. The model-based version is too small because the fitted inclusion surface only approximates the suppressed step selection, across which the exposure varies in mean and variance, so a covariance that treats the fit as exact under-covers.

The evaluation therefore reports coefficient error and interval behaviour separately. Empirical-standard-error oracle coverage replaces the reported standard error by the empirical spread of the estimates. Bias-corrected oracle coverage removes the average simulation bias while retaining each method's reported standard error. These are diagnostics based on hypothetical replacements. They show how coverage changes under each replacement but do not causally isolate a unique source of failure. Reporting them beside nominal coverage is more informative than a single coverage number \citep{morris2019using}.

\section{The generative methods under Design I}
\label{sec:conc_designI}

SICG and SIRD were also evaluated in Design~I, although neither uses the released design weight as a correction. SICG has exposure bias of $6.6\%$, an SE ratio of $0.26$, nominal coverage of $0.37$, and empirical-standard-error oracle coverage of $0.90$. Its reported standard errors are too small relative to the empirical spread. SIRD has exposure bias of $19.8\%$, an SE ratio of $1.50$, nominal coverage of $0.60$, empirical-standard-error oracle coverage of $0.21$, and bias-corrected oracle coverage of $0.98$. Its standard errors are conservative on average, while the oracle pattern points to displacement from the uncorrected informative selection.

Repeated refitting cannot replace a sampling correction. SIRD propagates more fitted-model variation than SICG, but it does not use the released weight to correct the informative selection. Weighted-likelihood multiple imputation and raking address that part of Design~I directly and have much lower exposure bias.

A structural difference between the two generative methods is worth naming, because it bears on why neither closes this gap. SIRD's path diffuses only the phase-two reference measurement toward its proxy, so its denoising objective, a variational bound that decomposes into Kullback--Leibler terms over the reference-measurement trajectory \citep{sohl-dickstein2015deep,ho2020denoising}, treats a conditioning covariate such as the design weight as an input to the reverse kernel rather than as part of the distribution it matches. SICG's critic instead concatenates its inputs, so the adversarial divergence \citep{goodfellow2014generative} is evaluated on the completed record together with those covariates, a more direct channel through which the weight can act. This difference, rather than the proxy anchoring, is the more likely reason SICG's Design~I exposure bias of $6.6\%$ is well below SIRD's $19.8\%$: the anchoring serves SIRD well in Designs~II and~III, where the measurement error is larger and the proxy is further from the reference measurement, so it is not what fails when the proxy is close, as in Design~I. Even so, neither channel is a design-based correction: both networks are trained on the selected rows and treat the weight as a feature rather than an importance weight, so neither reweights the likelihood or calibrates to the inclusion probabilities, and neither reproduces what weighted-likelihood imputation, raking and SISL achieve directly.

The same distinction appears on the Cox product term: a conservative reported standard error does not correct a displaced coefficient estimate.

\section{The two-phase setting and the gap between two literatures}
\label{sec:conc_gap}

Deep generative imputation is often evaluated by reconstruction, distributional similarity or downstream prediction \citep{yoon2018gain,tashiro2021csdi}. Statistical multiple imputation asks whether estimates and intervals from the completed datasets have suitable repeated-sampling behaviour \citep{Rubin1987,meng1994multiple,naf2023imputation}. Good performance under the first set of criteria does not establish the second. The chapters therefore tune network components with held-out criteria where needed, then retain or reject methods using coefficient bias and coverage.

GAIN and TabCSDI perform poorly on several inferential summaries in this thesis. In Design~II, GAIN has exposure bias above $50\%$ on the continuous block, and both methods report standard errors well below the empirical spread. The experiments do not isolate one cause. The training sample is restricted to the validated subsample, most reference values are missing, and the retained objectives were not designed around the target coefficients. SIRD combines a proxy-referenced denoising path with bootstrap refitting and gives lower coefficient error and coverage closer to $0.95$ than SICG on the main coefficient blocks. The comparison evaluates the full SIRD procedure and does not estimate the separate contribution of each component. Benchmark results from other missingness patterns do not verify inference in a two-phase validation design.

\section{Computational cost and outlook}
\label{sec:conc_cost}

Appendix~\ref{app:runtime} reports the full per-replicate cost. \texttt{mixgb} ranges from seconds to minutes. The chained-equations methods take tens of seconds in the 21-variable design and about two to three hours in the 119-variable Design~III. TabCSDI takes about eight to ten and a half hours per replicate in the runs reported here. SICG and SIRD take from minutes to about an hour. In Design~III, SIRD takes about $3{,}800$ seconds and SICG about $600$ seconds. This difference is largely the price of SIRD fitting 20 bootstrap denoisers while SICG trains once and uses dropout at imputation. The longer SIRD runtime is part of its refitting-based uncertainty procedure. Appendix~\ref{app:hyperparam} reports the comparator searches, Appendix~\ref{app:proposed_config} reports the retained deep-model configurations, and Chapters~\ref{chap:SICG} and~\ref{chap:SIRD} report the structural ablations.

Future hardware and software may reduce this cost, but the thesis makes no projection about how quickly that will occur. The present comparison should therefore be read at the reported runtime and hardware configuration. A cheaper procedure still needs repeated-sampling validation before it can replace the retained bootstrap.

\section{Limitations}
\label{sec:conc_limitations}

The main inferential limitation is the Cox-model age-by-insulin interaction. SICG and SIRD both have severe displacement and very low coverage on this term, and so do essentially all of the imputation comparators, while design-based raking stays close to nominal. That pattern points to uncongeniality \citep{meng1994multiple}: the completed datasets do not preserve the insulin-by-age relationship the analysis model uses, and the failure is shared by every imputation route regardless of architecture rather than produced by the hazard model itself, which raking fits without failing. Design~II gives a different interaction result: SIRD has aggregate nominal coverage of $0.94$, so the Cox failure should not be generalised to every product term.

SISL improves on naive imputation in Design~I but has more exposure bias than weighted-likelihood multiple imputation and raking. Its model-based version also underestimates the uncertainty, because the fitted inclusion surface only approximates the suppressed step selection. This conclusion is limited to the tested method and working models.

SICG gives reasonable categorical coefficient estimates but under-covers. SIRD reduces the categorical bias to $2.1\%$ and has nominal coverage of $0.99$, although its SE ratio of $1.40$ is conservative and the tree procedures have lower RMSE. SIRD also costs 20 fitted denoisers per replicate. VCCC is a single cohort and is used as an applied coefficient-estimate check rather than independent evidence of repeated-sampling calibration.

\section{Directions for further work}
\label{sec:conc_future}

The Cox-interaction failure is a congeniality problem, so the natural remedy is a congenial imputation that preserves the insulin-by-age interaction the analysis uses, either by including that interaction in the imputation model or by matching the completed-data cross-moment to the analysis. The present methods impute the exposure without imposing that structure, and whether building it in removes the displacement is untested. The categorical component could also be compared with alternatives to the transition-prior construction, such as a discrete analogue of the residual path.

The retained bootstrap injects fitted-model variation by training a separate denoiser on each resample, which is an ensemble of diffusion refits and is therefore expensive. A small but active line of recent work estimates the epistemic, or fitted-model, uncertainty of a diffusion model without a full ensemble, using hyper-networks and other single-model or low-cost constructions \citep{chan2024estimating,berry2024shedding,ekmekci2023quantifying,jazbec2025generative,gupta2026quantifying}, and this is the natural place to look for a cheaper substitute for the bootstrap. The obstacle is the target rather than the technique: these methods are built and benchmarked for the predictive uncertainty of imaging and forecasting tasks, and none are evaluated against Rubin's between-imputation variance under a two-phase analysis, so each would have to be tested one at a time for multiple imputation, judged by coefficient coverage and SE ratios over repeated replicates rather than by a predictive score.

Multimodality in the imputed continuous variables is a further open problem. Tabular GANs handle it through mode-specific normalisation, where a variational Gaussian mixture encodes each value as a mode indicator together with a within-mode value \citep{xu2019modeling,zhao2022ctabganplus}. In an imputation task this makes the objective harder, because the model then has to recover both the mode and the value from the small validated sample at once. SICG's residual construction could reduce the problem by tying each draw to the phase-one proxy, but it does so by shrinking the imputation spread, the under-dispersion seen earlier in this chapter. SIRD instead tends to return a unimodal density, because its reverse process is Gaussian at each step. The diffusion literature usually handles this with a quantile normalisation that maps each continuous margin to a Gaussian before diffusion \citep{kotelnikov2023tabddpm}, but that transform has to be estimated from the validated rows, and the five percent validated here is too small to fit it stably. A method that recovers a multimodal density and still returns a calibrated coefficient and interval is left for future work.

The tuning objective for the machine-learning imputers is also open. The distance-based criteria used here and in the wider literature, such as the energy distance \citep{szekely2003statistics,naf2025mar}, reward reproduction of the marginal or joint distribution but are only loosely tied to the bias and coverage of a downstream coefficient. A tuning target that is analysis-relevant yet not specialised to one analysis model would be more useful; no such standard is established, and finding one is a problem in its own right.

The generative methods also trained only on the validated phase-two subsample and used the phase-one information solely as conditioning. The much larger phase-one sample carries information about the covariate distribution and the proxy that could help training, and semi-supervised training is the natural device. Two obstacles are separate. The adversarial generator--discriminator training of the WGAN is a structural difficulty in its own right, because that training game sits awkwardly with the small validated sample of a two-phase design. The failure of semi-supervised training is broader and is not specific to the WGAN: the evaluation gap applies to any of these objectives, adversarial or not, none of which is aligned with the between-imputation variance or the target coefficient, so additional phase-one signal need not translate into better inference. The variants attempted here did not help, among them a teacher--student construction, in which a teacher network fitted on the validated rows supplies targets on the unlabelled phase-one records to train a student network \citep{tarvainen2017mean}. Whether a semi-supervised objective aligned with the inferential target can use phase-one information without degrading coverage is unresolved.

Further studies should also vary proxy strength, validation fraction, analysis nonlinearity and the number of imputed covariates. External information about a confidential sampling design is a further extension for SISL, provided it can be supplied without disclosing protected design variables.

\section{Concluding remarks}
\label{sec:conc_final}

The thesis studied automated imputation for two-phase validation with one target held fixed: an unbiased coefficient and a correctly covering interval. The results divide into two settings. When informative phase-two selection is represented only by a released design weight, as in Design~I, the weight has to enter the correction directly, and weighted-likelihood multiple imputation and generalised raking, and to a smaller degree SISL, do better than an imputation model that only conditions on the weight. When the selection-related phase-one variables are available, the task is the conditional distribution of the reference measurement, and a generative imputer can compete: SIRD has lower continuous and categorical bias than SICG and coverage closer to $0.95$ on those blocks, with Design~II continuous coverage $0.93$, Design~III continuous coverage $0.89$ under conservative standard errors, and categorical coverage $0.99$.

The two generative chapters make one diagnosis and one partial fix. SICG shows that a single fitted network can give low coefficient bias and still report intervals that are too narrow, because sampling its noise varies the imputations but not the fitted parameters. SIRD shows that the missing variation can be restored by refitting, which the phase-one proxy makes affordable by leaving the network only a residual to learn. Neither result removes the underlying gap: matching the data distribution, the criterion used in the machine-learning literature, does not guarantee a calibrated coefficient, and a method can do the first and fail the second.

Bootstrap refitting raises the between-imputation variance above the single-fit baseline but overshoots it rather than calibrating it; the Cox age-by-insulin interaction is badly covered by every generative method, so a completed dataset should not be treated as suitable for an arbitrary downstream term; and in Design~I neither generative method recovers what the released weight already carries. Generalised raking stays the lower-bias reference where its calibration constraints hold. The standard the thesis returns to is plain: an imputation method meant for inference must be judged by the target coefficient and the reported variance of that coefficient, because completed values that recover the data well from one fitted model do not, on their own, give a calibrated interval \citep{naf2023imputation,Rubin1987}.
